\chapter{Conclusiones y Trabajo Futuro}
\label{cap:Conclusiones}

Este capítulo ofrece una síntesis de la investigación realizada, valorando el grado de cumplimiento de los objetivos iniciales y destacando los hallazgos más relevantes obtenidos tras la extensa variedad de experimentos. Asimismo, se exponen las limitaciones encontradas, se proponen líneas de investigación futuras para dar continuidad al proyecto, y se detalla y justifica el esfuerzo temporal y computacional invertido en su desarrollo.

\section{Conclusiones}
\subsection*{Consecución de Objetivos}
La evaluación de los objetivos planteados al inicio de esta memoria muestra un grado de cumplimiento muy satisfactorio en términos generales, si bien es necesario realizar una valoración individualizada para señalar aquellos que se han alcanzado de forma parcial:
\begin{itemize}
    \item \textbf{Estudiar el estado del arte:} Cumplido. Se ha fundamentado la investigación en la literatura reciente sobre \textit{Deep Learning}, PLN y Visión Artificial aplicados a la detección de contenido de odio.
    \item \textbf{Preprocesar y adecuar el conjunto de datos:} Cumplido. Se desarrolló un \textit{pipeline} robusto de limpieza textual mediante \textit{regex} y extracción temporal de \textit{frames} mediante \textit{OpenCV}, además de implementar lógicas matemáticas para el manejo de etiquetas bajo el paradigma \textit{LeWiDi}.
    \item \textbf{Implementar y ajustar (\textit{fine-tuning}) modelos:} Cumplido. Se entrenaron y compararon con éxito arquitecturas discriminativas (\textit{RoBERTa}), generativas (\textit{Mistral}), de visión estática (\textit{ConvNeXt}) y de visión temporal (\textit{TimeSFormer}).
    \item \textbf{Proponer y desarrollar estrategias de fusión:} Cumplido. Se diseñó, optimizó y evaluó un sistema \textit{Ensemble} basado en \textit{Late Fusion} mediante la suma ponderada de probabilidades.
    \item \textbf{Participar en la tarea 3.1 de \textit{EXIST 2025}:} Cumplido. Se enviaron las predicciones oficiales utilizando \textit{BERT} y \textit{Llama} (\textit{zero-shot}), estableciendo el \textit{baseline} del proyecto.
    \item \textbf{Desarrollar marco de evaluación local (Tarea 3.1 y 3.3):} \textbf{Alcanzado parcialmente}. Si bien se generó un entorno de evaluación riguroso, el marco local de la Tarea 3.1 presentó un fuerte sesgo (\textit{overfitting}) al reutilizar el \textit{holdout} para la selección de hiperparámetros y la ponderación. Asimismo, la Tarea 3.3 no se abordó como un sistema multietiqueta jerárquico completo, sino que se condicionó a un análisis exploratorio simplificado (\textit{One-vs-Rest}) que sufrió un colapso predictivo ante el desbalanceo extremo de las clases minoritarias.
    \item \textbf{Explorar el impacto del perspectivismo:} Cumplido. Se demostró la sensibilidad algorítmica a la perspectiva de género mediante la segmentación demográfica de los anotadores y técnicas de \textit{Persona Prompting}.
\end{itemize}

\subsection*{Resumen de los Hallazgos}
Los resultados de la experimentación sugieren fuertemente que el sexismo en plataformas digitales es un fenómeno multimodal. Ningún modelo unimodal logró igualar el rendimiento del sistema unificado mediante \textit{Late Fusion} en el entorno de validación local (donde alcanzó un \textit{F1-score} de 0.7261). Al evaluar su capacidad predictiva frente al conjunto ciego oficial de la competición, el \textit{Ensemble} obtuvo un \textit{F1-score} de 0.6311. Esta caída de rendimiento evidencia la existencia de un sobreajuste (\textit{overfitting}) previo durante la fase de selección, pero, aun así, la arquitectura unificada logró una mejora objetiva frente a los enfoques individuales. El análisis exploratorio detallado indicó que el texto resulta especialmente útil para detectar cargas ideológicas, mientras que la información visual y temporal parece aportar un contexto valioso para abordar fenómenos complejos como la cosificación.

Además, los experimentos de perspectivismo han mostrado observaciones empíricas de gran valor para el diseño de IAs: la anotación del discurso de odio realizada desde una perspectiva femenina resultó en modelos con métricas de recuperación (\textit{Recall}) más sólidas. Por su parte, la simulación de roles en Grandes Modelos de Lenguaje demostró que inyectar el rol ``masculino'' en el \textit{prompt} dispara la tasa de falsos negativos. Si bien estos hallazgos están condicionados por el tamaño de los \textit{datasets} y los sesgos estadísticos preentrenados en los \textit{LLMs}, sugieren que la identidad demográfica (ya sea real o simulada) altera la evaluación algorítmica del discurso de odio.

\subsection*{Implicaciones}
Los resultados obtenidos respaldan el potencial de combinar \textit{LLMs} ajustados mediante \textit{PEFT} (\textit{QLoRA}) con \textit{Vision Transformers}. Se observa empíricamente que, para moderar plataformas modernas de vídeos cortos, el análisis de texto aislado puede resultar insuficiente, recomendándose su hibridación con la semántica del movimiento y la disposición espacial para comprender la verdadera intencionalidad del contenido.

\subsection*{Limitaciones}
La principal limitación técnica encontrada ha sido el desbalanceo extremo en las clases minoritarias del \textit{dataset} multietiqueta (especialmente en \textit{Sexual Violence} y \textit{Misogyny and Non-Sexual Violence}). Esta carencia de ejemplos positivos provocó un colapso predictivo en todas las familias de modelos, que optaron por predecir matemáticamente la clase mayoritaria. Por otro lado, hemos tenido la limitación derivada de haber reutilizado el mismo conjunto local de desarrollo para la selección de hiperparámetros y pesos del \textit{Ensemble}, lo que generó un sesgo al alza en las métricas internas reportadas. Finalmente, la extracción del texto incrustado en el vídeo introdujo ruido y ``basura algorítmica'' (como \textit{tags} y menciones) que dificultó la comprensión semántica de las redes.

\subsection*{Contribuciones}
Este trabajo aporta una metodología documentada de manera exhaustiva para el tratamiento multimodal de datos en español e inglés. Se ha analizado el comportamiento comparativo de modelos discriminativos y generativos, se han evaluado distintas lógicas estadísticas de agregación de fotogramas (\textit{Mean-Pooling, Majority-Vote}), y se ha explorado el impacto de los sesgos demográficos (\textit{LeWiDi}) en el entrenamiento de la inteligencia artificial.

\subsection*{Consideraciones Éticas}
El uso de modelos de inteligencia artificial para la moderación de contenido humano conlleva profundas implicaciones éticas. Como ha ilustrado el análisis de errores y de perspectivismo de este documento, los modelos tienden a heredar e incluso amplificar los sesgos de sus datos de entrenamiento. Las decisiones de censura o bloqueo automatizadas en redes sociales no deben ser tomadas como verdades absolutas, sino que deben contextualizarse y contar con supervisión humana para evitar la discriminación y el silenciamiento sistémico. 

\section{Trabajo Futuro}
\subsection*{Mejoras Metodológicas y Preprocesamiento}
A corto plazo, una línea de trabajo esencial consiste en implementar filtros de limpieza de texto mucho más robustos para mitigar el ruido introducido por el \textit{OCR} y las transcripciones de audio automáticas, traduciendo los metadatos visuales a un formato puramente semántico. Asimismo, resultaría crítico aplicar técnicas de \textit{Oversampling} avanzado o penalización de pesos funcionales (\textit{Class Weights}) para solventar el colapso predictivo de los modelos ante el desbalanceo de las categorías más graves.

\subsection*{Extensiones del Estudio Multimodal}
Aunque este trabajo ha fusionado texto, imagen estática y secuencias temporales de vídeo, el sonido aporta información vital no verbal (tono de voz, sarcasmo acústico, risas). Una extensión natural del sistema consistirá en integrar modelos especializados en características acústicas, como \textit{Wav2Vec} o \textit{Audio Spectogram Transformers}, añadiendo una nueva dimensión a la ecuación del \textit{Ensemble}. Además, sería interesante investigar técnicas de \textit{Early Fusion} o arquitecturas de atención cruzada (\textit{Cross-Attention}) que permitan a los modelos interconectar y evaluar texto e imagen de forma simultánea, y no en etapas separadas.


\subsection*{Aplicaciones Prácticas}
A nivel aplicado, la arquitectura híbrida documentada podría sentar las bases para una herramienta de \textit{software} capaz de asistir a los equipos de \textit{Trust \& Safety} de plataformas digitales. Esta herramienta no buscaría censurar autónomamente, sino priorizar y marcar de forma automática aquellos contenidos que, por su combinación de factores textuales y visuales, presenten una alta probabilidad de albergar discursos de odio, optimizando así el tiempo de revisión humana.

\section{Planificación Temporal del Trabajo Realizado}
En este apartado se detalla el tiempo y los recursos empleados en la investigación y redacción del presente Trabajo de Fin de Máster. Dada la magnitud de la experimentación requerida, el mayor esfuerzo no recayó en las entregas iniciales de la competición, sino en el análisis avanzado, entrenamiento y depuración de modelos a posteriori. Las horas globales estimadas se resumen en la \autoref{tab:PlanificationTemporal}.

\begin{table}[H]
    \centering
    \begin{tabularx}{\textwidth}{Xr}
        \toprule
        \textbf{Tarea} & \textbf{Horas Estimadas} \\
        \midrule
        \textbf{Revisión del Estado del Arte e Investigación Teórica:} & \textbf{50} \\
        \quad - Análisis del discurso de odio, \textit{Transformers}, y \textit{LLMs}. & \\
        \quad - Estudio del Perspectivismo (\textit{LeWiDi}) y métricas de evaluación. & \\
        \addlinespace[0.2em]
        \textbf{Familiarización con el Entorno y Tecnologías:} & \textbf{50} \\
        \quad - Despliegue en \textit{Google Colab}, librerías (\textit{PyTorch}, \textit{Hugging Face}). & \\
        \quad - Lógicas de extracción multimedia (\textit{OpenCV}, \textit{Decord}). & \\
        \addlinespace[0.2em]
        \textbf{Tratamiento de Datos y Entregas Preliminares (EXIST 2025):} & \textbf{70} \\
        \quad - Limpieza de \textit{datasets}, gestión de nulos y estructuración de \textit{frames}. & \\
        \quad - \textit{Splits} estratificados para la tarea binaria y multietiqueta. & \\
        \addlinespace[0.2em]
        \textbf{Experimentación Avanzada \textit{a posteriori} y Desarrollo:} & \textbf{140} \\
        \quad - Ejecución y entrenamiento intensivo en GPUs (entorno L4 de Colab). & \\
        \quad - \textit{Debugging}, resolución de errores de memoria (OOM), HPO. & \\
        \quad - Algoritmos de \textit{Late Fusion} y \textit{Persona Prompting} en \textit{Mistral}. & \\
        \addlinespace[0.2em]
        \textbf{Elaboración de la Memoria y Formato:} & \textbf{70} \\
        \quad - Análisis cruzado de resultados, extracción de gráficas y conclusiones. & \\
        \quad - Redacción intensiva, maquetación en \LaTeX{} y revisión final. & \\
        \midrule
        \textbf{Total} & \textbf{380} \\
        \bottomrule
    \end{tabularx}
    \caption{Estimación de la planificación temporal del trabajo realizado.}
    \label{tab:PlanificationTemporal}
\end{table}

Como se refleja en la tabla, el proyecto ha requerido una inversión aproximada de 380 horas. Cabe destacar el elevado coste computacional de la fase de Experimentación Avanzada. Durante esta etapa, se consumieron más de 300 unidades de cómputo en la plataforma \textit{Google Colab} utilizando hardware de alto rendimiento (\textit{GPUs NVIDIA L4}) para poder procesar los grandes volúmenes de datos de imágenes y vídeo, y realizar el \textit{Fine-Tuning} de arquitecturas pesadas como \textit{ConvNeXt}, \textit{TimeSFormer} y el modelo de lenguaje \textit{Mistral-7B}. Este volumen de experimentación es el que ha permitido estructurar y fundamentar los hallazgos descritos a lo largo de este documento. 